% Chương 1 — Giới thiệu. Được nạp bằng % Chương 1 — Giới thiệu. Được nạp bằng % Chương 1 — Giới thiệu. Được nạp bằng % Chương 1 — Giới thiệu. Được nạp bằng \input{chapter1} trong main.tex.
% Chỉ dùng các gói đã khai báo trong main.tex; không thêm gói mới.
\chapter{GIỚI THIỆU}

\section{Bối cảnh và lý do chọn đề tài}

Truy xuất văn bản là bài toán tìm trong một kho tài liệu những văn bản liên quan
nhất tới một câu hỏi của người dùng. Với tiếng Việt, bài toán này khó hơn so với
các ngôn ngữ dùng khoảng trắng để phân tách từ, vì khoảng trắng trong tiếng Việt
phân tách \emph{âm tiết} chứ không phân tách \emph{từ}: một từ như ``trí tuệ nhân
tạo'' gồm bốn âm tiết được viết rời nhau. Nếu lập chỉ mục trên từng âm tiết, các
âm tiết riêng lẻ như ``tuệ'' hay ``tạo'' trở thành đơn vị chỉ mục rời rạc và làm
loãng tín hiệu phân biệt của tài liệu. Vì vậy tách từ tiếng Việt là bước bắt buộc
trước khi biểu diễn và lập chỉ mục văn bản.

Một khó khăn thứ hai là lọc từ dừng. Loại các từ chức năng rất phổ biến (như
``của'', ``và'', ``là'') làm giảm nhiễu và kích thước chỉ mục, nhưng danh sách
cố định không tính đến đặc thù của từng miền dữ liệu. Trong kho tin tức tiếng
Việt theo nhiều chủ đề, có những từ phổ biến nhưng chỉ tập trung ở vài chủ đề;
loại chúng đi có thể làm mất thông tin phân biệt. Đề tài này tiếp cận bằng một
tầng lọc thứ hai dựa trên độ đo cấu trúc hình học chứ không dựa trên tần suất đặt
tay, nhằm trả lời câu hỏi: khi bỏ một từ khỏi không gian biểu diễn, các tài liệu
có trở nên dễ phân biệt nhau hơn hay không?

Đề tài được thực hiện trong phạm vi môn học đa phương tiện, với dữ liệu, thuật
toán và quy trình đo đều nằm lại trong kho mã nguồn của nhóm. Yêu cầu gốc của bài
tập nhóm bao gồm tách từ, loại từ dừng, biểu diễn TF-IDF, tính tương đồng cosin
và đánh giá kết quả \cite{baigiang}.

\section{Mục tiêu của đồ án}

Đồ án hướng tới các mục tiêu sau:

\begin{enumerate}
  \item Xây dựng một quy trình xử lý ngôn ngữ tiếng Việt hoàn chỉnh và có thể
        kiểm chứng: chuẩn hóa Unicode, tách từ theo ba phương pháp, và lọc từ
        dừng theo hai tầng.
  \item Cài đặt một hệ thống truy xuất văn bản trên nền TF-IDF và tương đồng
        cosin, có chỉ mục đảo, có khả năng giải thích điểm số theo từng từ.
  \item Cài đặt độ đo Giá trị phân biệt từ (Term Discrimination Value -- TDV)
        một cách chính xác để tự động phát hiện từ dừng đặc thù của miền dữ liệu.
  \item Bốn cấu hình thực nghiệm đối chứng (C1--C4) để so sánh các chiến lược
        lọc từ dừng trên cùng một kho ngữ liệu, cùng một bộ truy vấn và cùng một
        bộ nhãn liên quan.
  \item Đo độ trễ truy vấn và dung lượng chỉ mục bằng số liệu thật, có quy trình
        đo tách bạch khỏi các chỉ số chất lượng phụ thuộc nhãn người.
\end{enumerate}

\section{Phạm vi số liệu}

Kho ngữ liệu của đồ án gồm 2.500 bài báo tiếng Việt thuộc 10 lĩnh vực, mỗi lĩnh
vực 250 bài, được khóa bằng hàm băm SHA256 trong
\texttt{data/corpus\_manifest.json}. Bộ truy vấn gồm 20 câu, khoảng hai câu mỗi
lĩnh vực, phủ các dạng ngắn, dài, phủ định và chuyên ngành. Tập ứng viên đánh giá
được gộp từ kết quả của các đường cơ sở và được đưa cho người gán nhãn.

Một điểm được nêu rõ ngay từ đầu: báo cáo này \textbf{phân biệt rõ} giữa phần đã
đo và phần chưa đo. Trạng thái nhãn liên quan tại thời điểm viết là
\texttt{pending\_human\_annotation}: pool ứng viên đã hình thành nhưng \textbf{chưa
có cặp nào được người gán xác nhận}. Do đó các chỉ số chất lượng truy xuất
(P@K, R@K, F1@K, MAP, NDCG@K) \textbf{chưa được đo} và không được trình bày như
một kết quả đã có. Việc này là quyết định thiết kế: một bảng toàn số không khi
thiếu nhãn dễ bị đọc nhầm là hệ thống không tìm được gì, trong khi nguyên nhân là
thiếu dữ liệu gán nhãn.

\section{Yêu cầu của đề tài}

Theo yêu cầu bài tập nhóm \cite{baigiang}, sản phẩm cần thể hiện đủ ba phần lý
thuyết, cài đặt và thực nghiệm. Cụ thể đồ án phải có: (i) tách từ tiếng Việt bằng
\emph{so khớp cực đại} trên từ điển và cây Trie; (ii) loại từ dừng có nguồn và
quy tắc rõ ràng, không tự bù danh sách để đủ số lượng; (iii) biểu diễn văn bản và
truy vấn bằng TF-IDF, xếp hạng bằng tương đồng cosin; (iv) bộ công cụ đo lường và
biểu đồ kết quả. Nền tảng lý thuyết về biểu diễn văn bản và chỉ mục đảo được tham
chiếu từ giáo trình môn học \cite{giaotrinh}.

\section{Tổ chức thực hiện}

Công việc được chia thành ba nhóm trách nhiệm, ranh giới giữa các nhóm được ghi
thành hợp đồng API trong kho mã nguồn để tránh chồng lấn:

\begin{itemize}
  \item \textbf{M1 -- Xử lý ngôn ngữ tự nhiên:} chuẩn hóa chuỗi, từ điển và cây
        Trie, tách từ FMM/BMM/BiMM, lọc từ dừng tầng 1 cùng danh sách bảo vệ
        phủ định.
  \item \textbf{M2 -- Truy xuất thông tin:} TF-IDF, chỉ mục đảo, tương đồng
        cosin, độ đo TDV tầng 2, bộ đánh giá chỉ số và bộ đo độ trễ, dung lượng.
  \item \textbf{M3 -- Tích hợp và giao diện:} ứng dụng minh họa, hợp đồng API,
        quản trị mã nguồn và đóng gói.
\end{itemize}

Đồ án được thực hiện bởi Nhóm 7. Báo cáo dùng cách gọi nhất quán ``Nhóm 7'' trong
toàn bộ tài liệu.

\section{Đóng góp chính}

Các đóng góp được trình bày trong báo cáo gồm: một quy trình tách từ tiếng Việt
có vết thuật toán từng bước, tái lập được; một cài đặt TDV chính xác có đường
tham chiếu độc lập để kiểm chứng, kèm cả một phiên bản ``cách làm sai'' có chủ
đích để chứng minh rằng việc tái chuẩn hóa sau khi xóa một chiều là bắt buộc; một
bộ đánh giá có cổng chặn chống công bố kết quả giả khi thiếu nhãn hoặc sai hàm
băm; và một quy trình đo chất lượng tách từ bằng nhãn người với ràng buộc toàn vẹn
theo hàm băm.

\section{Cấu trúc báo cáo}

Phần còn lại của báo cáo được tổ chức như sau. Chương~2 trình bày nền tảng xử lý
ngôn ngữ tiếng Việt và lý do của từng bước tiền xử lý. Chương~3 trình bày thiết kế
hệ thống tìm kiếm, bao gồm thiết kế tầng ngôn ngữ M1, độ đo TF-IDF, độ đo TDV và
bốn cấu hình lọc. Chương~4 trình bày thực nghiệm và đánh giá. Chương~5 kết luận,
nêu rõ phần đã kiểm chứng và phần còn chờ nhãn người, cùng các hạn chế và hướng
phát triển.
 trong main.tex.
% Chỉ dùng các gói đã khai báo trong main.tex; không thêm gói mới.
\chapter{GIỚI THIỆU}

\section{Bối cảnh và lý do chọn đề tài}

Truy xuất văn bản là bài toán tìm trong một kho tài liệu những văn bản liên quan
nhất tới một câu hỏi của người dùng. Với tiếng Việt, bài toán này khó hơn so với
các ngôn ngữ dùng khoảng trắng để phân tách từ, vì khoảng trắng trong tiếng Việt
phân tách \emph{âm tiết} chứ không phân tách \emph{từ}: một từ như ``trí tuệ nhân
tạo'' gồm bốn âm tiết được viết rời nhau. Nếu lập chỉ mục trên từng âm tiết, các
âm tiết riêng lẻ như ``tuệ'' hay ``tạo'' trở thành đơn vị chỉ mục rời rạc và làm
loãng tín hiệu phân biệt của tài liệu. Vì vậy tách từ tiếng Việt là bước bắt buộc
trước khi biểu diễn và lập chỉ mục văn bản.

Một khó khăn thứ hai là lọc từ dừng. Loại các từ chức năng rất phổ biến (như
``của'', ``và'', ``là'') làm giảm nhiễu và kích thước chỉ mục, nhưng danh sách
cố định không tính đến đặc thù của từng miền dữ liệu. Trong kho tin tức tiếng
Việt theo nhiều chủ đề, có những từ phổ biến nhưng chỉ tập trung ở vài chủ đề;
loại chúng đi có thể làm mất thông tin phân biệt. Đề tài này tiếp cận bằng một
tầng lọc thứ hai dựa trên độ đo cấu trúc hình học chứ không dựa trên tần suất đặt
tay, nhằm trả lời câu hỏi: khi bỏ một từ khỏi không gian biểu diễn, các tài liệu
có trở nên dễ phân biệt nhau hơn hay không?

Đề tài được thực hiện trong phạm vi môn học đa phương tiện, với dữ liệu, thuật
toán và quy trình đo đều nằm lại trong kho mã nguồn của nhóm. Yêu cầu gốc của bài
tập nhóm bao gồm tách từ, loại từ dừng, biểu diễn TF-IDF, tính tương đồng cosin
và đánh giá kết quả \cite{baigiang}.

\section{Mục tiêu của đồ án}

Đồ án hướng tới các mục tiêu sau:

\begin{enumerate}
  \item Xây dựng một quy trình xử lý ngôn ngữ tiếng Việt hoàn chỉnh và có thể
        kiểm chứng: chuẩn hóa Unicode, tách từ theo ba phương pháp, và lọc từ
        dừng theo hai tầng.
  \item Cài đặt một hệ thống truy xuất văn bản trên nền TF-IDF và tương đồng
        cosin, có chỉ mục đảo, có khả năng giải thích điểm số theo từng từ.
  \item Cài đặt độ đo Giá trị phân biệt từ (Term Discrimination Value -- TDV)
        một cách chính xác để tự động phát hiện từ dừng đặc thù của miền dữ liệu.
  \item Bốn cấu hình thực nghiệm đối chứng (C1--C4) để so sánh các chiến lược
        lọc từ dừng trên cùng một kho ngữ liệu, cùng một bộ truy vấn và cùng một
        bộ nhãn liên quan.
  \item Đo độ trễ truy vấn và dung lượng chỉ mục bằng số liệu thật, có quy trình
        đo tách bạch khỏi các chỉ số chất lượng phụ thuộc nhãn người.
\end{enumerate}

\section{Phạm vi số liệu}

Kho ngữ liệu của đồ án gồm 2.500 bài báo tiếng Việt thuộc 10 lĩnh vực, mỗi lĩnh
vực 250 bài, được khóa bằng hàm băm SHA256 trong
\texttt{data/corpus\_manifest.json}. Bộ truy vấn gồm 20 câu, khoảng hai câu mỗi
lĩnh vực, phủ các dạng ngắn, dài, phủ định và chuyên ngành. Tập ứng viên đánh giá
được gộp từ kết quả của các đường cơ sở và được đưa cho người gán nhãn.

Một điểm được nêu rõ ngay từ đầu: báo cáo này \textbf{phân biệt rõ} giữa phần đã
đo và phần chưa đo. Trạng thái nhãn liên quan tại thời điểm viết là
\texttt{pending\_human\_annotation}: pool ứng viên đã hình thành nhưng \textbf{chưa
có cặp nào được người gán xác nhận}. Do đó các chỉ số chất lượng truy xuất
(P@K, R@K, F1@K, MAP, NDCG@K) \textbf{chưa được đo} và không được trình bày như
một kết quả đã có. Việc này là quyết định thiết kế: một bảng toàn số không khi
thiếu nhãn dễ bị đọc nhầm là hệ thống không tìm được gì, trong khi nguyên nhân là
thiếu dữ liệu gán nhãn.

\section{Yêu cầu của đề tài}

Theo yêu cầu bài tập nhóm \cite{baigiang}, sản phẩm cần thể hiện đủ ba phần lý
thuyết, cài đặt và thực nghiệm. Cụ thể đồ án phải có: (i) tách từ tiếng Việt bằng
\emph{so khớp cực đại} trên từ điển và cây Trie; (ii) loại từ dừng có nguồn và
quy tắc rõ ràng, không tự bù danh sách để đủ số lượng; (iii) biểu diễn văn bản và
truy vấn bằng TF-IDF, xếp hạng bằng tương đồng cosin; (iv) bộ công cụ đo lường và
biểu đồ kết quả. Nền tảng lý thuyết về biểu diễn văn bản và chỉ mục đảo được tham
chiếu từ giáo trình môn học \cite{giaotrinh}.

\section{Tổ chức thực hiện}

Công việc được chia thành ba nhóm trách nhiệm, ranh giới giữa các nhóm được ghi
thành hợp đồng API trong kho mã nguồn để tránh chồng lấn:

\begin{itemize}
  \item \textbf{M1 -- Xử lý ngôn ngữ tự nhiên:} chuẩn hóa chuỗi, từ điển và cây
        Trie, tách từ FMM/BMM/BiMM, lọc từ dừng tầng 1 cùng danh sách bảo vệ
        phủ định.
  \item \textbf{M2 -- Truy xuất thông tin:} TF-IDF, chỉ mục đảo, tương đồng
        cosin, độ đo TDV tầng 2, bộ đánh giá chỉ số và bộ đo độ trễ, dung lượng.
  \item \textbf{M3 -- Tích hợp và giao diện:} ứng dụng minh họa, hợp đồng API,
        quản trị mã nguồn và đóng gói.
\end{itemize}

Đồ án được thực hiện bởi Nhóm 7. Báo cáo dùng cách gọi nhất quán ``Nhóm 7'' trong
toàn bộ tài liệu.

\section{Đóng góp chính}

Các đóng góp được trình bày trong báo cáo gồm: một quy trình tách từ tiếng Việt
có vết thuật toán từng bước, tái lập được; một cài đặt TDV chính xác có đường
tham chiếu độc lập để kiểm chứng, kèm cả một phiên bản ``cách làm sai'' có chủ
đích để chứng minh rằng việc tái chuẩn hóa sau khi xóa một chiều là bắt buộc; một
bộ đánh giá có cổng chặn chống công bố kết quả giả khi thiếu nhãn hoặc sai hàm
băm; và một quy trình đo chất lượng tách từ bằng nhãn người với ràng buộc toàn vẹn
theo hàm băm.

\section{Cấu trúc báo cáo}

Phần còn lại của báo cáo được tổ chức như sau. Chương~2 trình bày nền tảng xử lý
ngôn ngữ tiếng Việt và lý do của từng bước tiền xử lý. Chương~3 trình bày thiết kế
hệ thống tìm kiếm, bao gồm thiết kế tầng ngôn ngữ M1, độ đo TF-IDF, độ đo TDV và
bốn cấu hình lọc. Chương~4 trình bày thực nghiệm và đánh giá. Chương~5 kết luận,
nêu rõ phần đã kiểm chứng và phần còn chờ nhãn người, cùng các hạn chế và hướng
phát triển.
 trong main.tex.
% Chỉ dùng các gói đã khai báo trong main.tex; không thêm gói mới.
\chapter{GIỚI THIỆU}

\section{Bối cảnh và lý do chọn đề tài}

Truy xuất văn bản là bài toán tìm trong một kho tài liệu những văn bản liên quan
nhất tới một câu hỏi của người dùng. Với tiếng Việt, bài toán này khó hơn so với
các ngôn ngữ dùng khoảng trắng để phân tách từ, vì khoảng trắng trong tiếng Việt
phân tách \emph{âm tiết} chứ không phân tách \emph{từ}: một từ như ``trí tuệ nhân
tạo'' gồm bốn âm tiết được viết rời nhau. Nếu lập chỉ mục trên từng âm tiết, các
âm tiết riêng lẻ như ``tuệ'' hay ``tạo'' trở thành đơn vị chỉ mục rời rạc và làm
loãng tín hiệu phân biệt của tài liệu. Vì vậy tách từ tiếng Việt là bước bắt buộc
trước khi biểu diễn và lập chỉ mục văn bản.

Một khó khăn thứ hai là lọc từ dừng. Loại các từ chức năng rất phổ biến (như
``của'', ``và'', ``là'') làm giảm nhiễu và kích thước chỉ mục, nhưng danh sách
cố định không tính đến đặc thù của từng miền dữ liệu. Trong kho tin tức tiếng
Việt theo nhiều chủ đề, có những từ phổ biến nhưng chỉ tập trung ở vài chủ đề;
loại chúng đi có thể làm mất thông tin phân biệt. Đề tài này tiếp cận bằng một
tầng lọc thứ hai dựa trên độ đo cấu trúc hình học chứ không dựa trên tần suất đặt
tay, nhằm trả lời câu hỏi: khi bỏ một từ khỏi không gian biểu diễn, các tài liệu
có trở nên dễ phân biệt nhau hơn hay không?

Đề tài được thực hiện trong phạm vi môn học đa phương tiện, với dữ liệu, thuật
toán và quy trình đo đều nằm lại trong kho mã nguồn của nhóm. Yêu cầu gốc của bài
tập nhóm bao gồm tách từ, loại từ dừng, biểu diễn TF-IDF, tính tương đồng cosin
và đánh giá kết quả \cite{baigiang}.

\section{Mục tiêu của đồ án}

Đồ án hướng tới các mục tiêu sau:

\begin{enumerate}
  \item Xây dựng một quy trình xử lý ngôn ngữ tiếng Việt hoàn chỉnh và có thể
        kiểm chứng: chuẩn hóa Unicode, tách từ theo ba phương pháp, và lọc từ
        dừng theo hai tầng.
  \item Cài đặt một hệ thống truy xuất văn bản trên nền TF-IDF và tương đồng
        cosin, có chỉ mục đảo, có khả năng giải thích điểm số theo từng từ.
  \item Cài đặt độ đo Giá trị phân biệt từ (Term Discrimination Value -- TDV)
        một cách chính xác để tự động phát hiện từ dừng đặc thù của miền dữ liệu.
  \item Bốn cấu hình thực nghiệm đối chứng (C1--C4) để so sánh các chiến lược
        lọc từ dừng trên cùng một kho ngữ liệu, cùng một bộ truy vấn và cùng một
        bộ nhãn liên quan.
  \item Đo độ trễ truy vấn và dung lượng chỉ mục bằng số liệu thật, có quy trình
        đo tách bạch khỏi các chỉ số chất lượng phụ thuộc nhãn người.
\end{enumerate}

\section{Phạm vi số liệu}

Kho ngữ liệu của đồ án gồm 2.500 bài báo tiếng Việt thuộc 10 lĩnh vực, mỗi lĩnh
vực 250 bài, được khóa bằng hàm băm SHA256 trong
\texttt{data/corpus\_manifest.json}. Bộ truy vấn gồm 20 câu, khoảng hai câu mỗi
lĩnh vực, phủ các dạng ngắn, dài, phủ định và chuyên ngành. Tập ứng viên đánh giá
được gộp từ kết quả của các đường cơ sở và được đưa cho người gán nhãn.

Một điểm được nêu rõ ngay từ đầu: báo cáo này \textbf{phân biệt rõ} giữa phần đã
đo và phần chưa đo. Trạng thái nhãn liên quan tại thời điểm viết là
\texttt{pending\_human\_annotation}: pool ứng viên đã hình thành nhưng \textbf{chưa
có cặp nào được người gán xác nhận}. Do đó các chỉ số chất lượng truy xuất
(P@K, R@K, F1@K, MAP, NDCG@K) \textbf{chưa được đo} và không được trình bày như
một kết quả đã có. Việc này là quyết định thiết kế: một bảng toàn số không khi
thiếu nhãn dễ bị đọc nhầm là hệ thống không tìm được gì, trong khi nguyên nhân là
thiếu dữ liệu gán nhãn.

\section{Yêu cầu của đề tài}

Theo yêu cầu bài tập nhóm \cite{baigiang}, sản phẩm cần thể hiện đủ ba phần lý
thuyết, cài đặt và thực nghiệm. Cụ thể đồ án phải có: (i) tách từ tiếng Việt bằng
\emph{so khớp cực đại} trên từ điển và cây Trie; (ii) loại từ dừng có nguồn và
quy tắc rõ ràng, không tự bù danh sách để đủ số lượng; (iii) biểu diễn văn bản và
truy vấn bằng TF-IDF, xếp hạng bằng tương đồng cosin; (iv) bộ công cụ đo lường và
biểu đồ kết quả. Nền tảng lý thuyết về biểu diễn văn bản và chỉ mục đảo được tham
chiếu từ giáo trình môn học \cite{giaotrinh}.

\section{Tổ chức thực hiện}

Công việc được chia thành ba nhóm trách nhiệm, ranh giới giữa các nhóm được ghi
thành hợp đồng API trong kho mã nguồn để tránh chồng lấn:

\begin{itemize}
  \item \textbf{M1 -- Xử lý ngôn ngữ tự nhiên:} chuẩn hóa chuỗi, từ điển và cây
        Trie, tách từ FMM/BMM/BiMM, lọc từ dừng tầng 1 cùng danh sách bảo vệ
        phủ định.
  \item \textbf{M2 -- Truy xuất thông tin:} TF-IDF, chỉ mục đảo, tương đồng
        cosin, độ đo TDV tầng 2, bộ đánh giá chỉ số và bộ đo độ trễ, dung lượng.
  \item \textbf{M3 -- Tích hợp và giao diện:} ứng dụng minh họa, hợp đồng API,
        quản trị mã nguồn và đóng gói.
\end{itemize}

Đồ án được thực hiện bởi Nhóm 7. Báo cáo dùng cách gọi nhất quán ``Nhóm 7'' trong
toàn bộ tài liệu.

\section{Đóng góp chính}

Các đóng góp được trình bày trong báo cáo gồm: một quy trình tách từ tiếng Việt
có vết thuật toán từng bước, tái lập được; một cài đặt TDV chính xác có đường
tham chiếu độc lập để kiểm chứng, kèm cả một phiên bản ``cách làm sai'' có chủ
đích để chứng minh rằng việc tái chuẩn hóa sau khi xóa một chiều là bắt buộc; một
bộ đánh giá có cổng chặn chống công bố kết quả giả khi thiếu nhãn hoặc sai hàm
băm; và một quy trình đo chất lượng tách từ bằng nhãn người với ràng buộc toàn vẹn
theo hàm băm.

\section{Cấu trúc báo cáo}

Phần còn lại của báo cáo được tổ chức như sau. Chương~2 trình bày nền tảng xử lý
ngôn ngữ tiếng Việt và lý do của từng bước tiền xử lý. Chương~3 trình bày thiết kế
hệ thống tìm kiếm, bao gồm thiết kế tầng ngôn ngữ M1, độ đo TF-IDF, độ đo TDV và
bốn cấu hình lọc. Chương~4 trình bày thực nghiệm và đánh giá. Chương~5 kết luận,
nêu rõ phần đã kiểm chứng và phần còn chờ nhãn người, cùng các hạn chế và hướng
phát triển.
 trong main.tex.
% Chỉ dùng các gói đã khai báo trong main.tex; không thêm gói mới.
\chapter{GIỚI THIỆU}

\section{Bối cảnh và lý do chọn đề tài}

Truy xuất văn bản là bài toán tìm trong một kho tài liệu những văn bản liên quan
nhất tới một câu hỏi của người dùng. Với tiếng Việt, bài toán này khó hơn so với
các ngôn ngữ dùng khoảng trắng để phân tách từ, vì khoảng trắng trong tiếng Việt
phân tách \emph{âm tiết} chứ không phân tách \emph{từ}: một từ như ``trí tuệ nhân
tạo'' gồm bốn âm tiết được viết rời nhau. Nếu lập chỉ mục trên từng âm tiết, các
âm tiết riêng lẻ như ``tuệ'' hay ``tạo'' trở thành đơn vị chỉ mục rời rạc và làm
loãng tín hiệu phân biệt của tài liệu. Vì vậy tách từ tiếng Việt là bước bắt buộc
trước khi biểu diễn và lập chỉ mục văn bản.

Một khó khăn thứ hai là lọc từ dừng. Loại các từ chức năng rất phổ biến (như
``của'', ``và'', ``là'') làm giảm nhiễu và kích thước chỉ mục, nhưng danh sách
cố định không tính đến đặc thù của từng miền dữ liệu. Trong kho tin tức tiếng
Việt theo nhiều chủ đề, có những từ phổ biến nhưng chỉ tập trung ở vài chủ đề;
loại chúng đi có thể làm mất thông tin phân biệt. Đề tài này tiếp cận bằng một
tầng lọc thứ hai dựa trên độ đo cấu trúc hình học chứ không dựa trên tần suất đặt
tay, nhằm trả lời câu hỏi: khi bỏ một từ khỏi không gian biểu diễn, các tài liệu
có trở nên dễ phân biệt nhau hơn hay không?

Đề tài được thực hiện trong phạm vi môn học đa phương tiện, với dữ liệu, thuật
toán và quy trình đo đều nằm lại trong kho mã nguồn của nhóm. Yêu cầu gốc của bài
tập nhóm bao gồm tách từ, loại từ dừng, biểu diễn TF-IDF, tính tương đồng cosin
và đánh giá kết quả \cite{baigiang}.

\section{Mục tiêu của đồ án}

Đồ án hướng tới các mục tiêu sau:

\begin{enumerate}
  \item Xây dựng một quy trình xử lý ngôn ngữ tiếng Việt hoàn chỉnh và có thể
        kiểm chứng: chuẩn hóa Unicode, tách từ theo ba phương pháp, và lọc từ
        dừng theo hai tầng.
  \item Cài đặt một hệ thống truy xuất văn bản trên nền TF-IDF và tương đồng
        cosin, có chỉ mục đảo, có khả năng giải thích điểm số theo từng từ.
  \item Cài đặt độ đo Giá trị phân biệt từ (Term Discrimination Value -- TDV)
        một cách chính xác để tự động phát hiện từ dừng đặc thù của miền dữ liệu.
  \item Bốn cấu hình thực nghiệm đối chứng (C1--C4) để so sánh các chiến lược
        lọc từ dừng trên cùng một kho ngữ liệu, cùng một bộ truy vấn và cùng một
        bộ nhãn liên quan.
  \item Đo độ trễ truy vấn và dung lượng chỉ mục bằng số liệu thật, có quy trình
        đo tách bạch khỏi các chỉ số chất lượng phụ thuộc nhãn người.
\end{enumerate}

\section{Phạm vi số liệu}

Kho ngữ liệu của đồ án gồm 2.500 bài báo tiếng Việt thuộc 10 lĩnh vực, mỗi lĩnh
vực 250 bài, được khóa bằng hàm băm SHA256 trong
\texttt{data/corpus\_manifest.json}. Bộ truy vấn gồm 20 câu, khoảng hai câu mỗi
lĩnh vực, phủ các dạng ngắn, dài, phủ định và chuyên ngành. Tập ứng viên đánh giá
được gộp từ kết quả của các đường cơ sở và được đưa cho người gán nhãn.

Một điểm được nêu rõ ngay từ đầu: báo cáo này \textbf{phân biệt rõ} giữa phần đã
đo và phần chưa đo. Trạng thái nhãn liên quan tại thời điểm viết là
\texttt{pending\_human\_annotation}: pool ứng viên đã hình thành nhưng \textbf{chưa
có cặp nào được người gán xác nhận}. Do đó các chỉ số chất lượng truy xuất
(P@K, R@K, F1@K, MAP, NDCG@K) \textbf{chưa được đo} và không được trình bày như
một kết quả đã có. Việc này là quyết định thiết kế: một bảng toàn số không khi
thiếu nhãn dễ bị đọc nhầm là hệ thống không tìm được gì, trong khi nguyên nhân là
thiếu dữ liệu gán nhãn.

\section{Yêu cầu của đề tài}

Theo yêu cầu bài tập nhóm \cite{baigiang}, sản phẩm cần thể hiện đủ ba phần lý
thuyết, cài đặt và thực nghiệm. Cụ thể đồ án phải có: (i) tách từ tiếng Việt bằng
\emph{so khớp cực đại} trên từ điển và cây Trie; (ii) loại từ dừng có nguồn và
quy tắc rõ ràng, không tự bù danh sách để đủ số lượng; (iii) biểu diễn văn bản và
truy vấn bằng TF-IDF, xếp hạng bằng tương đồng cosin; (iv) bộ công cụ đo lường và
biểu đồ kết quả. Nền tảng lý thuyết về biểu diễn văn bản và chỉ mục đảo được tham
chiếu từ giáo trình môn học \cite{giaotrinh}.

\section{Tổ chức thực hiện}

Công việc được chia thành ba nhóm trách nhiệm, ranh giới giữa các nhóm được ghi
thành hợp đồng API trong kho mã nguồn để tránh chồng lấn:

\begin{itemize}
  \item \textbf{M1 -- Xử lý ngôn ngữ tự nhiên:} chuẩn hóa chuỗi, từ điển và cây
        Trie, tách từ FMM/BMM/BiMM, lọc từ dừng tầng 1 cùng danh sách bảo vệ
        phủ định.
  \item \textbf{M2 -- Truy xuất thông tin:} TF-IDF, chỉ mục đảo, tương đồng
        cosin, độ đo TDV tầng 2, bộ đánh giá chỉ số và bộ đo độ trễ, dung lượng.
  \item \textbf{M3 -- Tích hợp và giao diện:} ứng dụng minh họa, hợp đồng API,
        quản trị mã nguồn và đóng gói.
\end{itemize}

Đồ án được thực hiện bởi Nhóm 7. Báo cáo dùng cách gọi nhất quán ``Nhóm 7'' trong
toàn bộ tài liệu.

\section{Đóng góp chính}

Các đóng góp được trình bày trong báo cáo gồm: một quy trình tách từ tiếng Việt
có vết thuật toán từng bước, tái lập được; một cài đặt TDV chính xác có đường
tham chiếu độc lập để kiểm chứng, kèm cả một phiên bản ``cách làm sai'' có chủ
đích để chứng minh rằng việc tái chuẩn hóa sau khi xóa một chiều là bắt buộc; một
bộ đánh giá có cổng chặn chống công bố kết quả giả khi thiếu nhãn hoặc sai hàm
băm; và một quy trình đo chất lượng tách từ bằng nhãn người với ràng buộc toàn vẹn
theo hàm băm.

\section{Cấu trúc báo cáo}

Phần còn lại của báo cáo được tổ chức như sau. Chương~2 trình bày nền tảng xử lý
ngôn ngữ tiếng Việt và lý do của từng bước tiền xử lý. Chương~3 trình bày thiết kế
hệ thống tìm kiếm, bao gồm thiết kế tầng ngôn ngữ M1, độ đo TF-IDF, độ đo TDV và
bốn cấu hình lọc. Chương~4 trình bày thực nghiệm và đánh giá. Chương~5 kết luận,
nêu rõ phần đã kiểm chứng và phần còn chờ nhãn người, cùng các hạn chế và hướng
phát triển.
